% Macros for Paper B.
\newtheorem{theorem}{Theorem}[section]
\newtheorem{lemma}[theorem]{Lemma}
\newtheorem{proposition}[theorem]{Proposition}
\newtheorem{corollary}[theorem]{Corollary}
\theoremstyle{definition}
\newtheorem{definition}[theorem]{Definition}
\theoremstyle{remark}
\newtheorem{remark}[theorem]{Remark}
\renewcommand{\Re}{\operatorname{Re}}
\renewcommand{\Im}{\operatorname{Im}}
\newcommand{\tr}{\operatorname{tr}}
\newcommand{\abs}[1]{\lvert #1\rvert}
\newcommand{\norm}[1]{\lVert #1\rVert}
\newcommand{\floor}[1]{\lfloor #1\rfloor}
\newcommand{\R}{\mathbb R}
\newcommand{\C}{\mathbb C}
\newcommand{\Z}{\mathbb Z}
\newcommand{\eps}{\varepsilon}
\newcommand{\FQ}{\mathfrak F_Q}
\newcommand{\Farey}{\mathcal F}
\newcommand{\Lc}{\mathcal L}
\newcommand{\XX}{X}
\newcommand{\Lam}{\Lambda}
\newcommand{\DT}{D_T}
\newcommand{\Csieve}{C}
\newcommand{\alphap}{\alpha'}
\newcommand{\alphaz}{\alpha_0}
\newcommand{\Bfun}{\mathcal B}
\newcommand{\Om}{\Omega}
\newcommand{\Hw}{H_w}
\newcommand{\gwin}{g}
\newcommand{\Nsimp}{N^{s}_{0,\chi}}
\newcommand{\Nchi}{N_\chi}
\newcommand{\frob}[1]{\lVert #1\rVert_F}
% repository address and tag of the release (one place each)
\newcommand{\repourl}{\url{https://github.com/Oliverds321/simple-zeros}}
\newcommand{\repotag}{\texttt{v1.0}}
% the commit whose Lean sources the release contains
\newcommand{\relcommit}{\texttt{4874cc0}}
% the census quoted in sec. 1.5 and Appendix A (files audit/followup/census_2026-10-03_1557.* of the Lean repository);
% gen_appA_r5.py reads this definition and takes the census files of the commit that it names
\newcommand{\censusref}{the release commit \texttt{4874cc0}}
% Text that the author still has to confirm. To switch the marker off: \newcommand{\draftnote}[1]{#1}
\newcommand{\draftnote}[1]{#1\textbf{ [DRAFT, author to confirm]}}
% trust tags
\newcommand{\tagL}{\textup{[L]}}
\newcommand{\tagLm}{\textup{[L$\circ$]}}
\newcommand{\tagLs}{\textup{[L-stated]}}
\newcommand{\tagC}{\textup{[C]}}
\newcommand{\tagN}{\textup{[N]}}
\newcommand{\Vsim}{\textup{[V$\sim$]}}
\newcommand{\tagM}{\textup{[M]}}
